\section{Thiết kế nghiên cứu và dữ liệu}
\label{sec:method-design}

Nghiên cứu khảo sát các pipeline học máy và Ensemble Learning cho phát hiện xâm nhập nhị phân trên UNSW-NB15 \cite{moustafa2015unsw}. Nhãn Normal được mã hóa 0 và Attack là 1. Hai partition chính thức được giữ nguyên: TRAIN gồm 175\,341 dòng (56\,000 Normal; 119\,341 Attack), TEST gồm 82\,332 dòng (37\,000 Normal; 45\,332 Attack). Các trường \texttt{id}, \texttt{label} và \texttt{attack\_cat} không thuộc tập predictor. Trường \texttt{attack\_cat} chỉ được dùng sau dự đoán để mô tả các mẫu Attack bị bỏ sót; vì thế kết quả theo họ tấn công không phải phân loại đa lớp.

Thiết kế hướng tới so sánh một baseline Random Forest với Stacking, một cấu hình Random Forest có ít đặc trưng hơn, đồng thời khảo sát lỗi, độ bất định và chi phí quan sát được. Mục tiêu không phải đề xuất thuật toán ensemble mới. Novel-vector view được định nghĩa bằng việc loại khỏi TEST những dòng có vector 42 đặc trưng trùng với vector TRAIN; đây là phân tích độ nhạy trong cùng benchmark, không phải một split độc lập mới.

\section{Lựa chọn đặc trưng và tiền xử lý}
\label{sec:method-feature-selection}

Track 42 giữ toàn bộ đặc trưng dự báo sau khi loại các trường không dùng làm predictor. Track stable17 được xác định trên TRAIN bằng ba nhóm phương pháp filter, L1 và ExtraTrees. Với năm seed \texttt{11, 23, 42, 71, 101}, một biến được xem là ổn định trong từng phương pháp nếu xuất hiện ở ít nhất ba trên năm lượt; tập cuối giữ các biến đáp ứng ngưỡng ở tối thiểu hai phương pháp. Quy tắc này mô tả độ lặp lại của lựa chọn biến trong các lượt đã chạy, không phải bằng chứng về độ ổn định của hiệu năng dự đoán.

Thứ tự 17 biến được khóa là:
\begin{quote}\small\raggedright
\texttt{ct\_srv\_dst}, \texttt{ct\_state\_ttl}, \texttt{dload}, \texttt{dttl},
\texttt{proto}, \texttt{state}, \texttt{sttl}, \texttt{ct\_dst\_ltm},
\texttt{dmean}, \texttt{dwin}, \texttt{is\_sm\_ips\_ports},
\texttt{rate}, \texttt{service}, \texttt{swin},
\texttt{ct\_dst\_sport\_ltm}, \texttt{ct\_dst\_src\_ltm}, \texttt{dbytes}.
\end{quote}
Kết quả RFECV 13 biến là một track riêng, không thay cho stable17. Các biến phân loại \texttt{proto}, \texttt{service} và \texttt{state} được one-hot encode với \texttt{handle\_unknown=ignore}; biến số được chuẩn hóa trong pipeline Logistic Regression, còn Random Forest và Decision Tree dùng phép tiền xử lý phù hợp với estimator. Các phép biến đổi có tham số học được đặt trong pipeline để fit trên phần huấn luyện tương ứng.

Danh sách stable17 được tạo từ toàn bộ official TRAIN trước development cross-validation. Vì vậy, nhãn của các fold validation đã góp phần vào bước chọn biến và điểm development có thể lạc quan so với quy trình nested hoàn toàn. TEST không tham gia lựa chọn biến; tuy nhiên, sự phân biệt này không loại bỏ nguy cơ selection optimism trong cross-validation.

\section{Cấu trúc các pipeline}
\label{sec:method-pipelines}

Random Forest huấn luyện nhiều cây quyết định trên các mẫu bootstrap và tập đặc trưng ngẫu nhiên, sau đó tổng hợp xác suất hoặc phiếu dự đoán của các cây. Cấu hình RF-42 có 300 cây, \texttt{max\_features}=0.5 và \texttt{min\_samples\_leaf}=1; RF-17 có 150 cây với cùng tỷ lệ đặc trưng nhưng \texttt{min\_samples\_leaf}=2. Do nhiều tham số thay đổi cùng lúc, đối chiếu hai cấu hình không phải phép ablation chỉ thay số đặc trưng.

Voting kết hợp đầu ra của Logistic Regression (LR), Decision Tree (DT) và Random Forest (RF). Hard Voting chọn lớp theo đa số phiếu; Soft Voting lấy trung bình xác suất dự đoán trước khi áp dụng ngưỡng. Stacking cũng dùng ba learner này ở tầng cơ sở, nhưng huấn luyện Logistic Regression ở tầng kết hợp từ dự đoán ngoài fold của các learner cơ sở. Cách tạo đầu vào OOF làm giảm việc meta learner học trực tiếp từ dự đoán in-sample của tầng cơ sở; nó không tự loại bỏ mọi nguồn lạc quan trong quá trình chọn cấu hình.

\begin{table}[htbp]
\centering
\small
\renewcommand{\arraystretch}{1.3}
\setlength{\tabcolsep}{6pt}
\caption{Cấu hình của ba pipeline được dùng trong đánh giá cuối.}
\label{tab:frozen-candidates}
\begin{tabularx}{\linewidth}{>{\raggedright\arraybackslash}p{2.5cm}X S[table-format=2.0]}
\toprule
Pipeline & Cấu hình chính & {Số biến gốc}\\
\midrule
RF-42 & Random Forest: 300 cây, \texttt{max\_features}=0.5, \texttt{min\_samples\_leaf}=1. & 42\\
Stacking-42 & LR cơ sở: $C=1$, solver \texttt{liblinear}; DT: lá tối thiểu 2; RF như RF-42; meta LR: $C=0.5$, \texttt{liblinear}, \texttt{class\_weight=None}. Inner CV ba fold, \texttt{passthrough=False}. & 42\\
RF-17 & Random Forest: 150 cây, \texttt{max\_features}=0.5, \texttt{min\_samples\_leaf}=2. & 17\\
\bottomrule
\end{tabularx}
\end{table}

\section{Phát triển mô hình, OOF và nguyên tắc đánh giá}
\label{sec:method-validation}

Development OOF được tổng hợp từ outer \texttt{StratifiedKFold} ba fold với seed 42. Trong Stacking, inner \texttt{StratifiedKFold} ba fold dùng seed 137 để tạo dự đoán cho meta learner. Các kết quả Voting, Stacking, diversity và RF đối chiếu trong giai đoạn này chỉ phục vụ mô tả/phát triển mô hình trên TRAIN. Do stable17 được chọn trước outer CV và meta tuning chưa hoàn toàn nested, các con số OOF không được xem như ước lượng độc lập không thiên lệch.

Sau khi khóa ba cấu hình cuối, mỗi pipeline được fit lại trên toàn bộ TRAIN và dự đoán trên official TEST với ngưỡng 0.5. TEST không được dùng để đổi mô hình, threshold hay tập đặc trưng. Các chỉ số gồm confusion matrix, Precision, Recall, Macro-F1, MCC, FPR và FNR. Trường \texttt{pr\_auc} trong mã nguồn được tính bằng \texttt{average\_precision\_score}; vì vậy ở đây được gọi là Average Precision (AP), không đồng nhất với mọi cách tính diện tích dưới đường precision--recall.

Hiệu số giữa hai mô hình được đánh giá bằng paired bootstrap 10\,000 lượt với seed \texttt{20260928}. Các mẫu hình dự đoán của hai mô hình được lấy cùng nhau theo từng dòng, phân tầng theo nhãn thật, giữ nguyên tổng số lượng mỗi lớp. Khoảng 95\% là percentile CI riêng cho từng chỉ số, không điều chỉnh đồng thời toàn bộ các khoảng. McNemar exact so sánh tính đúng tổng thể của từng cặp dự đoán; Holm được áp dụng cho ba cặp mô hình trong mỗi view đánh giá. Phép thử này không kiểm định trực tiếp hiệu số Macro-F1. Cả bootstrap lẫn McNemar trên các dòng benchmark không ước lượng biến thiên giữa các lần huấn luyện độc lập.

\section{Khả năng tái lập trong môi trường phần mềm xác định}
\label{sec:method-p15}

Trong giai đoạn kiểm tra tính tái lập (P15), RF-17, RF-42 và Stacking-42 được tái dựng trên 175\,341 dòng official TRAIN trong môi trường Python 3.11.9 và scikit-learn 1.9.0, sau khi cấu hình và kết quả lịch sử đã được khóa. Kiểm tra artifact xác nhận hash, schema, khả năng nạp lại mô hình và tạo xác suất dự đoán trên các dòng TRAIN. Kết quả này cho thấy các cấu hình có thể được huấn luyện và sử dụng trong môi trường phần mềm được ghi nhận.

Môi trường và đối tượng mô hình lịch sử không được lưu đủ để chứng minh đồng nhất bitwise với các mô hình từng tạo dự đoán official TEST. Do đó, P15 là tái dựng/tái huấn luyện theo cấu hình đã ghi, không phải khôi phục chính xác artifact lịch sử và không phải phép xác minh lại các metric TEST. Kiểm tra predictability trên dữ liệu TRAIN-derived cũng không đánh giá khả năng vận hành của hệ thống IDS trong thực tế. Ba phạm vi này — khả năng tái lập trong môi trường xác định, tái huấn luyện từ TRAIN và triển khai — được xem là các tuyên bố riêng biệt.
